\subsection{Transformada Cuántica de Fourier: Variación de Fase y Baja Redundancia}
\label{subsec:patrones_qft}

La Transformada Cuántica de Fourier (QFT) aplicada a un estado base $\ket{k}$ produce una superposición con magnitudes uniformes pero fases dependientes del índice. En su forma ideal, la amplitud de cada componente $\ket{x}$ puede expresarse como:
\[
a_x = \frac{1}{\sqrt{N}} e^{\frac{2\pi i k x}{N}},
\]
lo cual implica que, aunque la magnitud sea constante, las partes real e imaginaria varían de forma oscilatoria con $x$.

Desde la perspectiva de almacenamiento, este comportamiento tiende a generar arreglos \textit{densos} donde la mayoría de entradas son no nulas y muchas de ellas difieren entre sí. En consecuencia, la redundancia directa (por repetición exacta) suele ser baja para compresión estrictamente sin pérdida. Esta intuición coincide con el tipo de deterioro del ratio reportado por \textcite{wu2018amplitude} cuando la evolución del circuito hace más diverso el contenido del vector. Por ello, la QFT constituye un caso razonable para esperar baja compresibilidad al final del circuito, aunque el comportamiento exacto depende del estado de entrada y no debe asumirse idéntico para toda variante del algoritmo.
